\chapter{書き込みと読み出し}
% full version: chapters/09_acet.tex

脳の記憶には不便な特徴がある。記憶を保存しているのは、仕事に使われているのと同じ結合なのだ。別のハードディスクはない。ネットワークは思い出すたびに結合を使い、学ぶたびに結合を変える。新しいことを書き込みながら、古いものをどう壊さずにおくのか。コンピュータには、そのための専用信号「書き込み許可」がある。脳では、その役を放送局\nt{ACET}、つまりアセチルコリンが担っているらしい。

\section*{2つのモード}

いくつかの絵を記憶しているネットワークを想像してほしい。半分を見せると、残りを描き足す。これが読み出しだ。次に、古い絵に似た新しい絵を見せる。ネットワークが癖で描き足してしまえば、新しいものではなく古いものを見てしまい、誤りを学ぶことになる。新しいことを書き込むには、ネットワークはしばらく自分の記憶を信じるのをやめ、世界を見なければならない。

\pencilfig{9}{\pencilart{9}}{1本の配線がネットワークを切り替える。新しいことを書き込むか、古いことを思い出すか。}

アセチルコリンはまさにこれを行い、その3つの効果はすべて測定されている。ネットワークが「思い出す」ための内部結合を弱め、感覚器から来る信号を強め、結合を変えやすくする。アセチルコリンが多いと書き込みモード。ネットワークは世界に耳を傾けて学ぶ。少ないと読み出しモード。ネットワークは記憶に頼って描き足す。

私たちのモデルには、両方がうまくいくようなレベルは存在しない。書き込みには内部結合を弱める必要があり、読み出しにはそれが最大の強さで必要だ。だから切り替えスイッチが欠かせない。

\section*{目をどれだけ信じるか}

エンジニアは似た問題をとうに解いている。船が航路を進むとき、航海士には船がいるはずの位置の推測航法による計算と、新しいが不正確な観測がある。それぞれをどれだけ信じるべきか。最適フィルタの答えはこうだ。計算に自信がないほど、観測を信じよ。そして同じ処方にはもう一つの教えがある。自信がないほど、速く学び直さなければならない。1つのつまみが入力の重みと学習の速さを同時に制御する。これはまさにアセチルコリンがしていることだ。ここから仮説が生まれる。アセチルコリンは、ネットワーク自身の世界モデルがどれだけ不確かかを伝えているのだ。

\pencilfig{124}{%
% optimal blending: the less sure the model, the more weight on the observation (and the faster it relearns)
\foreach \dx/\wm/\we/\ttop/\bbot in {0/2.4pt/0.6pt/{モデルは確信}/{アセチルコリン少}, 4.3/0.6pt/2.4pt/{モデルは不確か}/{アセチルコリン多}}{
  \begin{scope}[xshift=\dx cm]
  \draw[penlite=pcViolet, wash=pcViolet, rounded corners=3pt] (0,0.95) rectangle (1.3,1.45); \node[lbl, text=black] at (0.65,1.2) {記憶};
  \draw[penlite=pcBlue, wash=pcBlue, rounded corners=3pt] (0,-0.25) rectangle (1.3,0.25); \node[lbl, text=black] at (0.65,0) {目};
  \draw[dotpen=pcAmber, wash=pcAmber] (2.95,0.6) circle (0.55); \node[lbl, text=black] at (2.95,0.6) {推定};
  \draw[pcViolet, line width=\wm, line cap=round, -{Stealth[length=5pt]}] (1.35,1.2) -- (2.4,0.8);
  \draw[pcBlue, line width=\we, line cap=round, -{Stealth[length=5pt]}] (1.35,0) -- (2.4,0.4);
  \node[font=\small\itshape, text=pcGray, anchor=south] at (1.55,1.6) {\ttop};
  \node[lbl, anchor=north] at (1.55,-0.4) {\bbot};
  \end{scope}}
}{目をどれだけ信じるか。ネットワークが自分の世界モデルに自信があれば、推定はほぼすべて記憶から来る。自信がなければ観測から来て、そのときネットワークは速く学び直す。仮説では、このつまみを回しているのがアセチルコリンである。}

\section*{分業}

前の章を思い出そう。世界が突然変わると、仮説ではノルアドレナリンが予想外の出来事に気づく。そしてアセチルコリンは、新しい状況を学び終えるまでネットワークを書き込みモードに保つ。一方は「何かが変わった」と知らせ、もう一方は「書け」と命じる。

\section*{夜の書き直し}

深い睡眠の間、アセチルコリンのレベルは低い。ネットワークは世界から切り離され、読み出しモードにある。まさにこの時間に、昼間の出来事がネットワークの中で「再生」される。これは測定されている。その際に昼間学んだことが長期記憶に書き直されるというのは有力な仮説であり、睡眠の章で再び取り上げる。

\fullversion{9}
